% !TEX program = xelatex
\documentclass[12pt,a4paper]{ctexbook}

% ==================== 页面 ====================
\usepackage[top=2.5cm,bottom=2.5cm,left=2.5cm,right=2.5cm]{geometry}

% ==================== 字体 ====================
\usepackage{fontspec}
\usepackage{iffont}
\IfFontExistsTF{Consolas}
  {\setmonofont{Consolas}}
  {\setmonofont{DejaVu Sans Mono}[Scale=MatchLowercase]}

% ==================== 颜色 ====================
\usepackage{xcolor}
\definecolor{codebg}{HTML}{F5F5F5}
\definecolor{codeframe}{HTML}{E0E0E0}
\definecolor{cppkeyword}{HTML}{1A1AFF}
\definecolor{cppcomment}{HTML}{2E8B2E}
\definecolor{cppstring}{HTML}{B22222}
\definecolor{linenumgray}{HTML}{999999}
\definecolor{algheadbg}{HTML}{1A3C5E}
\definecolor{csharpkeyword}{HTML}{8B008B}
\definecolor{csharptype}{HTML}{006400}

% ==================== listings ====================
\usepackage{listings}

% ---- listing style ----
\lstdefinestyle{cppstyle}{
  language=C++,
  basicstyle=\small\ttfamily,
  keywordstyle=\color{cppkeyword}\bfseries,
  commentstyle=\color{cppcomment}\itshape,
  stringstyle=\color{cppstring},
  numberstyle=\tiny\color{linenumgray},
  numbers=left,
  frame=single,
  rulecolor=\color{codeframe},
  framesep=6pt,
  breaklines=true,
  tabsize=4,
  columns=flexible,
  keepspaces=true,
  backgroundcolor=\color{codebg},
  showstringspaces=false,
  morekeywords={constexpr,consteval,constinit,decltype,auto,
    concept,requires,co_await,co_yield,co_return,
    noexcept,override,final,nullptr,static_assert,
    template,typename,namespace,using,mutable,volatile,
    explicit,operator,friend,inline,virtual,
    thread_local,alignas,alignof,if,consteval,constexpr},
  deletekeywords={int,char,bool,float,double,long,short,unsigned,signed,void,wchar_t},
  emph={size_t,int32_t,uint32_t,int64_t,uint64_t,ssize_t,
    string,vector,map,set,unique_ptr,shared_ptr,optional,variant,any,
    span,string_view,
    ifstream,ofstream,fstream,iostream,
    ostream,istream,stringstream,ostringstream,
    path,filesystem},
  emphstyle=\color{teal!80!black},
  escapeinside={(*@}{@*)},
}

% ---- 文档特定语言定义 ----
\lstdefinelanguage{Cpp23}{%
  language=[GNU]C++,
  morekeywords={concept,requires,co_await,co_yield,co_return,consteval,constinit,this},
  morekeywords=[2]{std},
}
\lstset{style=cppstyle}

% ==================== tcolorbox ====================
\usepackage[most]{tcolorbox}

\newtcolorbox{guideline}[1][]{
  enhanced,breakable,
  colback=blue!5!white,colframe=blue!75!black,
  fonttitle=\bfseries,
  title=要点,#1,
  boxed title style={colback=blue!75!black},
  attach boxed title to top left={yshift=-2mm,xshift=2mm},
  arc=2mm,boxrule=0.8mm,
}
\newtcolorbox{compare}[1][]{
  enhanced,breakable,
  colback=green!3!white,colframe=green!50!black,
  fonttitle=\bfseries,
  title=对比,#1,
  boxed title style={colback=green!50!black},
  attach boxed title to top left={yshift=-2mm,xshift=2mm},
  arc=2mm,boxrule=0.8mm,
}
\newtcolorbox{warning}[1][]{
  enhanced,breakable,
  colback=red!5!white,colframe=red!75!black,
  fonttitle=\bfseries,
  title=常见陷阱,#1,
  boxed title style={colback=red!75!black},
  attach boxed title to top left={yshift=-2mm,xshift=2mm},
  arc=2mm,boxrule=0.8mm,
}
\newtcolorbox{note}[1][]{
  enhanced,breakable,
  colback=orange!5!white,colframe=orange!80!black,
  fonttitle=\bfseries,
  title=注意,#1,
  boxed title style={colback=orange!80!black},
  attach boxed title to top left={yshift=-2mm,xshift=2mm},
  arc=2mm,boxrule=0.8mm,
}
\newtcolorbox{bestpractice}[1][]{
  enhanced,breakable,
  colback=purple!5!white,colframe=purple!60!black,
  fonttitle=\bfseries,
  title=最佳实践,#1,
  boxed title style={colback=purple!60!black},
  attach boxed title to top left={yshift=-2mm,xshift=2mm},
  arc=2mm,boxrule=0.8mm,
}

% ==================== 章节标题 ====================
\usepackage{titlesec}
\usepackage{fancyhdr}
\usepackage{graphicx}
\usepackage{amssymb}
\usepackage{amsmath}
\usepackage{tabularx}
\usepackage{booktabs}
\usepackage{longtable}
\usepackage{array}
\usepackage{multirow}
\usepackage{enumitem}
\usepackage{seqsplit}

\titleformat{\chapter}[display]
  {\normalfont\huge\bfseries\color{algheadbg}}
  {\chaptertitlename\ \thechapter}{20pt}{\Huge}
\titleformat{\section}
  {\normalfont\Large\bfseries\color{blue!70!black}}
  {\thesection}{1em}{}
\titleformat{\subsection}
  {\normalfont\large\bfseries\color{blue!60!black}}
  {\thesubsection}{1em}{}

\pagestyle{fancy}
\fancyhf{}
\fancyhead[LE,RO]{\thepage}
\fancyhead[RE]{\leftmark}
\fancyhead[LO]{\rightmark}

\setlist[itemize]{leftmargin=2em,itemsep=2pt}
\setlist[enumerate]{leftmargin=2em,itemsep=2pt}

% ==================== 超链接 ====================
\usepackage{hyperref}
\hypersetup{
  colorlinks=true,
  linkcolor=blue!70!black,
  urlcolor=blue!60!black,
  bookmarks=true,
  pdfauthor={CRTP & PImpl Reference},
  pdftitle={C++ CRTP 与 PImpl 惯用法},
}

% ==================== 自定义命令 ====================
\newcommand{\cpp}[1]{C++#1}
\newcommand{\Fn}[1]{\texttt{#1()}}
\newcommand{\Tp}[1]{\texttt{#1}}

% ====================================================================
\begin{document}
\maketitle
\tableofcontents
\newpage

% ══════════════════════════════════════════════════════════════
\part{CRTP — 奇异递归模板模式}
% ══════════════════════════════════════════════════════════════

\section{概述}

CRTP（Curiously Recurring Template Pattern，奇异递归模板模式）是一种将派生类自身作为模板参数传递给基类的技术。它在\textbf{编译期}实现静态多态，避免了虚函数表的运行时开销，广泛应用于以下场景：

\begin{itemize}[leftmargin=2em]
  \item 静态多态（Static Polymorphism）：编译期接口分发
  \item 混入（Mixin）：为派生类注入通用功能
  \item 策略注入：通过基类模板参数组合行为
  \item 替代虚函数：在性能敏感代码中消除 vtable 开销
\end{itemize}

\section{传统 CRTP：无显式 \texttt{this}}

\subsection{核心机制}

传统 CRTP 的关键在于基类通过 \texttt{static\_cast} 将 \texttt{this} 指针转换为派生类指针，从而在编译期调用派生类的成员函数：

\begin{lstlisting}[caption={传统 CRTP 基本结构}]
template <typename Derived>
class Base {
public:
    // 通过 static_cast 向下转型
    void interface() {
        // 调用派生类的实现
        static_cast<Derived*>(this)->implementation();
    }

    // const 版本
    void interface() const {
        static_cast<const Derived*>(this)->implementation();
    }
};

class Concrete : public Base<Concrete> {
public:
    void implementation() {
        // 具体实现
    }
};
\end{lstlisting}

\subsection{完整案例：可计数混入}

以下是一个实用的对象计数器混入，每次构造和析构都会自动更新计数：

\begin{lstlisting}[caption={CRTP 计数器混入}]
template <typename T>
class Countable {
    static int count_;
public:
    Countable()  { ++count_; }
    ~Countable() { --count_; }

    // 只读接口
    static int getCount() { return count_; }

    // 禁止拷贝（防止计数混乱）
    Countable(const Countable&)            = delete;
    Countable& operator=(const Countable&) = delete;
};

// 每个类型拥有独立的静态计数
template <typename T>
int Countable<T>::count_ = 0;

// ---- 使用 ----
class Widget : public Countable<Widget> {
    std::string name_;
public:
    explicit Widget(std::string n) : name_(std::move(n)) {}
};

class Gadget : public Countable<Gadget> {};

int main() {
    Widget w1("A"), w2("B"), w3("C");
    Gadget g1, g2;

    std::cout << "Widgets: " << Widget::getCount()  // 3
              << ", Gadgets: " << Gadget::getCount() // 2
              << '\n';
    // w3 析构后：
    // Widgets: 2
}
\end{lstlisting}

\subsection{完整案例：可克隆接口}

利用 CRTP 实现深拷贝接口，替代虚函数 \texttt{clone()}：

\begin{lstlisting}[caption={CRTP 可克隆混入}]
template <typename Derived>
class Cloneable {
public:
    // 返回基类指针的多态克隆
    std::unique_ptr<Cloneable> clone() const {
        auto p = static_cast<const Derived*>(this);
        return std::make_unique<Derived>(*p);
    }

    // 返回具体类型的克隆（更方便）
    std::unique_ptr<Derived> cloneExact() const {
        auto p = static_cast<const Derived*>(this);
        return std::make_unique<Derived>(*p);
    }

    virtual ~Cloneable() = default;
};

class Shape : public Cloneable<Shape> {
    double area_;
public:
    explicit Shape(double a) : area_(a) {}
    double area() const { return area_; }
};

class Circle : public Cloneable<Circle> {
    double radius_;
public:
    explicit Circle(double r) : radius_(r) {}
    double area() const { return 3.14159 * radius_ * radius_; }
};
\end{lstlisting}

\subsection{完整案例：二元运算符自动生成}

CRTP 的经典用法——从一个运算符自动推导其余运算符：

\begin{lstlisting}[caption={CRTP 运算符混入}]
template <typename Derived>
class Comparable {
public:
    // 只需派生类实现 operator<
    friend bool operator>(const Derived& a, const Derived& b) {
        return b < a;
    }
    friend bool operator<=(const Derived& a, const Derived& b) {
        return !(b < a);
    }
    friend bool operator>=(const Derived& a, const Derived& b) {
        return !(a < b);
    }
    friend bool operator==(const Derived& a, const Derived& b) {
        return !(a < b) && !(b < a);
    }
    friend bool operator!=(const Derived& a, const Derived& b) {
        return !(a == b);
    }
};

struct Date : Comparable<Date> {
    int year, month, day;

    friend bool operator<(const Date& a, const Date& b) {
        if (a.year  != b.year)  return a.year  < b.year;
        if (a.month != b.month) return a.month < b.month;
        return a.day < b.day;
    }
};

// 使用：只需实现 <，自动获得 >, <=, >=, ==, !=
\end{lstlisting}

\subsection{传统 CRTP 的局限}

\begin{itemize}[leftmargin=2em]
  \item \textbf{\texttt{static\_cast} 不安全}：若派生类未正确继承 \texttt{Base<Derived>}，转换将产生未定义行为
  \item \textbf{代码可读性差}：\texttt{static\_cast<Derived*>(this)} 遍布基类代码
  \item \textbf{无法在 CRTP 基类中使用 \texttt{virtual}}：静态多态与动态多态混用时需要额外设计
  \item \textbf{调试困难}：深层模板实例化的错误信息难以阅读
\end{itemize}

% --------------------------------------------------------------------------------------------
\section{C++23 显式对象参数（Deducing \texttt{this}）}
% --------------------------------------------------------------------------------------------

\subsection{概述}

C++23 引入了\textbf{显式对象参数}（Explicit Object Parameter），也称为 ``Deducing this''。它允许在成员函数参数列表中显式声明 \texttt{this} 对象，从而通过模板推导自动获得正确的类型，完全消除 \texttt{static\_cast}：

\begin{lstlisting}[caption={C++23 显式对象参数语法}]
struct Widget {
    // "self" 是显式对象参数
    // 编译器自动推导其类型为 Widget（或 Widget&、const Widget& 等）
    void method(this Widget& self, int x) {
        self.data_ = x;       // 等价于 this->data_ = x
    }

    void method(this const Widget& self) {
        std::cout << self.data_;
    }

private:
    int data_ = 0;
};
\end{lstlisting}

\subsection{CRTP 的 C++23 重写}

使用显式对象参数，CRTP 基类不再需要模板参数，也不再需要 \texttt{static\_cast}：

\begin{lstlisting}[caption={C++23 无模板 CRTP}]
class Base {
public:
    // self 自动推导为实际派生类类型
    void interface(this auto& self) {
        self.implementation();    // 直接调用，无需转型
    }

    void interface(this const auto& self) {
        self.implementation();
    }

    // 返回值也可以自动推导
    auto getValue(this const auto& self) {
        return self.compute();
    }
};

class Concrete : public Base {
public:
    void implementation() {
        std::cout << "Concrete::implementation\n";
    }
    int compute() const { return 42; }
};

class Another : public Base {
public:
    void implementation() {
        std::cout << "Another::implementation\n";
    }
    int compute() const { return 99; }
};

int main() {
    Concrete c;
    c.interface();     // 输出: Concrete::implementation
    std::cout << c.getValue() << '\n';  // 42

    Another a;
    a.interface();     // 输出: Another::implementation
    std::cout << a.getValue() << '\n';  // 99
}
\end{lstlisting}

\subsection{案例重写：计数器混入（C++23 版）}

\begin{lstlisting}[caption={C++23 计数器混入——无模板}]
class Countable {
    int count_ = 0;
public:
    void onConstruct(this Countable& self) { ++self.count_; }
    void onDestruct(this Countable& self)  { --self.count_; }
    int  getCount(this const Countable& self) { return self.count_; }
};

// 使用：组合而非继承
class Widget {
    Countable counter_;
    std::string name_;
public:
    Widget(std::string n) : name_(std::move(n)) {
        counter_.onConstruct();
    }
    ~Widget() {
        counter_.onDestruct();
    }
    int getCount() const { return counter_.getCount(); }
};
\end{lstlisting}

\subsection{案例重写：运算符混入（C++23 版）}

\begin{lstlisting}[caption={C++23 运算符混入}]
class Comparable {
public:
    // 从 operator< 自动推导其余比较运算符
    friend bool operator>(this const auto& a, const auto& b) {
        return b < a;
    }
    friend bool operator<=(this const auto& a, const auto& b) {
        return !(b < a);
    }
    friend bool operator>=(this const auto& a, const auto& b) {
        return !(a < b);
    }
    friend bool operator==(this const auto& a, const auto& b) {
        return !(a < b) && !(b < a);
    }
    friend bool operator!=(this const auto& a, const auto& b) {
        return !(a == b);
    }
};

struct Version : Comparable {
    int major, minor, patch;

    friend bool operator<(const Version& a, const Version& b) {
        if (a.major != b.major) return a.major < b.major;
        if (a.minor != b.minor) return a.minor < b.minor;
        return a.patch < b.patch;
    }
};
\end{lstlisting}

\subsection{显式对象参数的额外优势}

\subsubsection{消除 cv-ref 重载样板}

传统写法需要为每种 cv-ref 限定符编写重载，显式对象参数可以用一个模板函数统一处理：

\begin{lstlisting}[caption={消除 cv-ref 重载}]
// ---- 传统写法：4 个重载 ----
struct Old {
    void process() &;
    void process() const&;
    void process() &&;
    void process() const&&;
};

// ---- C++23：1 个函数 ----
struct New {
    void process(this auto&& self) {
        // self 的类型完美保留了 cv-ref 限定
        // 可以配合 if constexpr 做分发
        if constexpr (std::is_const_v<std::remove_reference_t<decltype(self)>>) {
            // const 路径
        } else {
            // 非 const 路径
        }
    }
};
\end{lstlisting}

\subsubsection{安全递归 Lambda}

\begin{lstlisting}[caption={C++23 递归 Lambda}]
// 传统递归 Lambda 需要 std::function 或 Y 组合子
// C++23 可以直接在 lambda 中使用显式 this 参数

auto fib = [](this auto self, int n) -> int {
    if (n <= 1) return n;
    return self(n - 1) + self(n - 2);  // self 即 lambda 自身
};

std::cout << fib(10) << '\n';  // 55
\end{lstlisting}

\subsection{传统 CRTP 与 C++23 方式对比}

\begin{table}[h]
\centering
\caption{传统 CRTP vs C++23 Deducing this}
\label{tab:crtp-compare}
\begin{tabular}{@{}lll@{}}
\toprule
\textbf{特性} & \textbf{传统 CRTP} & \textbf{C++23 显式对象参数} \\
\midrule
基类是否需要模板  & 是 (\texttt{Base<Derived>})  & 否 \\
向下转型          & \texttt{static\_cast}        & 不需要 \\
类型安全          & 弱（错误继承 = UB）         & 强（编译器推导） \\
编译错误可读性    & 差（深层模板错误）          & 好 \\
cv-ref 重载       & 每个需单独写               & 一个 \texttt{auto\&\&} 统一 \\
标准版本          & C++98 起可用               & C++23 起可用 \\
递归 Lambda       & 需 \texttt{std::function}   & 原生支持 \\
\bottomrule
\end{tabular}
\end{table}

\newpage

% ══════════════════════════════════════════════════════════════
\part{PImpl — 指向实现的指针}
% ══════════════════════════════════════════════════════════════

\section{概述}

PImpl（Pointer to Implementation，指向实现的指针）是一种将类的公共接口与其内部实现解耦的惯用法。其核心思想是在公共类中仅保留一个指向私有实现类的指针，从而：

\begin{itemize}[leftmargin=2em]
  \item \textbf{编译防火墙}：头文件中不包含实现细节的依赖，减少编译时间
  \item \textbf{ABI 稳定性}：实现类的变更不影响公共类的二进制布局
  \item \textbf{头文件洁净}：用户头文件不暴露内部类型和前向声明
\end{itemize}

\section{裸指针版本}

\subsection{头文件（\texttt{widget.h}）}

裸指针版本需要手动管理内存，并严格遵守三五法则：

\begin{lstlisting}[caption={PImpl 裸指针版——头文件}]
// widget.h
#pragma once
#include <string>

class Widget {
public:
    // 构造函数和析构函数必须在 .cpp 中定义
    // 因为 Impl 在头文件中是不完整类型
    Widget(int id, std::string name);
    ~Widget();

    // 拷贝构造 / 赋值
    Widget(const Widget& other);
    Widget& operator=(const Widget& other);

    // 移动构造 / 赋值
    Widget(Widget&& other) noexcept;
    Widget& operator=(Widget&& other) noexcept;

    // 公共接口
    void doSomething();
    int  getId() const;
    const std::string& getName() const;

private:
    struct Impl;    // 前向声明，不完整类型
    Impl* pImpl_;   // 裸指针指向实现
};
\end{lstlisting}

\subsection{实现文件（\texttt{widget.cpp}）}

\begin{lstlisting}[caption={PImpl 裸指针版——实现文件}]
// widget.cpp
#include "widget.h"
#include <vector>
#include <iostream>

// ---- 完整定义（仅在 .cpp 中可见）----
struct Widget::Impl {
    int id_;
    std::string name_;
    std::vector<double> data_;

    Impl(int id, std::string name)
        : id_(id), name_(std::move(name)), data_(1000, 0.0) {}

    void doWork() {
        std::cout << "Widget[" << id_ << "] working...\n";
    }
};

// ---- 构造 / 析构 ----
Widget::Widget(int id, std::string name)
    : pImpl_(new Impl(id, std::move(name))) {}

Widget::~Widget() {
    delete pImpl_;
}

// ---- 拷贝 ----
Widget::Widget(const Widget& other)
    : pImpl_(new Impl(*other.pImpl_)) {}

Widget& Widget::operator=(const Widget& other) {
    if (this != &other) {
        // 先创建新对象，再删除旧对象（异常安全）
        Impl* temp = new Impl(*other.pImpl_);
        delete pImpl_;
        pImpl_ = temp;
    }
    return *this;
}

// ---- 移动 ----
Widget::Widget(Widget&& other) noexcept
    : pImpl_(other.pImpl_) {
    other.pImpl_ = nullptr;
}

Widget& Widget::operator=(Widget&& other) noexcept {
    if (this != &other) {
        delete pImpl_;
        pImpl_ = other.pImpl_;
        other.pImpl_ = nullptr;
    }
    return *this;
}

// ---- 公共接口转发 ----
void Widget::doSomething() {
    pImpl_->doWork();
}

int Widget::getId() const {
    return pImpl_->id_;
}

const std::string& Widget::getName() const {
    return pImpl_->name_;
}
\end{lstlisting}

\subsection{裸指针版本的注意事项}

\begin{itemize}[leftmargin=2em]
  \item 必须实现析构函数、拷贝构造/赋值、移动构造/赋值（规则五）
  \item 拷贝赋值中应采用 ``先构造后删除'' 策略以保证异常安全
  \item 移动后源对象的 \texttt{pImpl\_} 置为 \texttt{nullptr}，防止双重释放
  \item 所有使用 \texttt{Impl} 成员的函数定义必须放在 \texttt{.cpp} 文件中
\end{itemize}

% --------------------------------------------------------------------------------------------
\section{\texttt{unique\_ptr} 版本}
% --------------------------------------------------------------------------------------------

\subsection{头文件（\texttt{widget.h}）}

使用 \texttt{std::unique\_ptr} 自动管理生命周期，大幅减少手动内存管理代码：

\begin{lstlisting}[caption={PImpl unique\_ptr 版——头文件}]
// widget.h
#pragma once
#include <memory>
#include <string>

class Widget {
public:
    Widget(int id, std::string name);

    // 关键：即使使用 unique_ptr，析构函数仍必须在 .cpp 中定义
    // 因为 unique_ptr 的析构需要看到 Impl 的完整定义
    ~Widget();

    // 拷贝（unique_ptr 不可拷贝，需手动实现）
    Widget(const Widget& other);
    Widget& operator=(const Widget& other);

    // 移动（unique_ptr 可移动，但需要完整类型才能移动）
    Widget(Widget&& other) noexcept;
    Widget& operator=(Widget&& other) noexcept;

    // 公共接口
    void doSomething();
    int  getId() const;
    const std::string& getName() const;

private:
    struct Impl;
    std::unique_ptr<Impl> pImpl_;
};
\end{lstlisting}

\subsection{实现文件（\texttt{widget.cpp}）}

\begin{lstlisting}[caption={PImpl unique\_ptr 版——实现文件}]
// widget.cpp
#include "widget.h"
#include <vector>
#include <iostream>

// ---- 完整定义 ----
struct Widget::Impl {
    int id_;
    std::string name_;
    std::vector<double> data_;

    Impl(int id, std::string name)
        : id_(id), name_(std::move(name)), data_(1000, 0.0) {}

    void doWork() {
        std::cout << "Widget[" << id_ << "] working...\n";
    }
};

// ---- 构造 / 析构 ----
Widget::Widget(int id, std::string name)
    : pImpl_(std::make_unique<Impl>(id, std::move(name))) {}

Widget::~Widget() = default;   // 在此处 Impl 是完整类型

// ---- 拷贝 ----
Widget::Widget(const Widget& other)
    : pImpl_(std::make_unique<Impl>(*other.pImpl_)) {}

Widget& Widget::operator=(const Widget& other) {
    if (this != &other) {
        // unique_ptr 赋值：直接 reset 即可（异常安全）
        pImpl_ = std::make_unique<Impl>(*other.pImpl_);
    }
    return *this;
}

// ---- 移动 ----
Widget::Widget(Widget&& other) noexcept = default;
Widget& Widget::operator=(Widget&& other) noexcept = default;

// ---- 公共接口转发 ----
void Widget::doSomething() {
    pImpl_->doWork();
}

int Widget::getId() const {
    return pImpl_->id_;
}

const std::string& Widget::getName() const {
    return pImpl_->name_;
}
\end{lstlisting}

\subsection{\texttt{unique\_ptr} 版本的关键要点}

\subsubsection{为什么析构必须在 \texttt{.cpp} 中定义？}

即使使用了 \texttt{= default}，析构函数也必须在 \texttt{.cpp} 文件中显式声明为默认。原因是 \texttt{unique\_ptr<Impl>} 在析构时需要调用 \texttt{delete}，这要求 \texttt{Impl} 是完整类型。若析构在头文件中隐式生成，编译器在实例化析构时 \texttt{Impl} 仍为不完整类型，将导致编译错误：

\begin{lstlisting}[caption={常见错误：不完整类型的析构}]
// 错误！若 ~Widget() 在头文件中隐式生成
// unique_ptr<Impl> 的析构函数会在头文件中被实例化
// 此时 Impl 是不完整类型，编译报错：
//   "deletion of pointer to incomplete type"

class Widget {
    struct Impl;
    std::unique_ptr<Impl> pImpl_;
    // ~Widget() 隐式生成 -- 编译错误！
};
\end{lstlisting}

\subsubsection{移动操作为何也需要在 \texttt{.cpp} 中？}

\texttt{unique\_ptr} 的移动构造/赋值同样需要销毁目标对象的原有内容（调用 \texttt{delete}），因此也需要 \texttt{Impl} 是完整类型。将移动操作声明为 \texttt{= default} 并定义在 \texttt{.cpp} 中是正确的做法。

\newpage

% --------------------------------------------------------------------------------------------
\section{裸指针与 \texttt{unique\_ptr} 版本对比}
% --------------------------------------------------------------------------------------------

\begin{table}[h]
\centering
\caption{PImpl 裸指针 vs unique\_ptr}
\label{tab:pimpl-compare}
\begin{tabular}{@{}lll@{}}
\toprule
\textbf{特性} & \textbf{裸指针} & \textbf{\texttt{unique\_ptr}} \\
\midrule
内存管理        & 手动 \texttt{new/delete}      & 自动 RAII \\
析构函数        & 必须手动 \texttt{delete}      & \texttt{= default}（在 .cpp 中） \\
拷贝赋值异常安全 & 需 ``先构造后删除'' 技巧      & 直接 \texttt{make\_unique} \\
移动操作        & 手动置 \texttt{nullptr}       & \texttt{= default}（在 .cpp 中） \\
代码量          & 较多                          & 较少 \\
泄漏风险        & 高（异常路径可能泄漏）        & 无 \\
性能            & 无差异（编译器优化后相同）    & 无差异 \\
推荐度          & 遗留代码维护                  & 现代 C++ 首选 \\
\bottomrule
\end{tabular}
\end{table}

\section{进阶技巧}

\subsection{使用 \texttt{make\_unique} 的工厂方法}

将 PImpl 的构造封装为工厂函数，进一步隔离实现：

\begin{lstlisting}[caption={工厂方法封装 PImpl}]
// widget.h
class Widget {
public:
    static Widget create(int id, std::string name);
    // ...
};

// widget.cpp
Widget Widget::create(int id, std::string name) {
    return Widget(id, std::move(name));
}
\end{lstlisting}

\subsection{与 CRTP 结合：静态多态 PImpl}

将 CRTP 与 PImpl 结合，可以在编译期选择不同的实现策略：

\begin{lstlisting}[caption={CRTP + PImpl 结合}]
template <typename Impl>
class Handle {
    std::unique_ptr<Impl> pImpl_;
public:
    template <typename... Args>
    explicit Handle(Args&&... args)
        : pImpl_(std::make_unique<Impl>(std::forward<Args>(args)...)) {}

    // 通过 CRTP 风格转发
    auto perform() { return static_cast<Impl*>(pImpl_.get())->perform(); }
};

struct FastImpl {
    int perform() { /* 快速实现 */ return 1; }
};

struct SafeImpl {
    int perform() { /* 安全实现 */ return 2; }
};

// 编译期选择实现
using FastWidget = Handle<FastImpl>;
using SafeWidget = Handle<SafeImpl>;
\end{lstlisting}

\subsection{PImpl 与 \texttt{shared\_ptr}}

当需要共享所有权语义时，可以使用 \texttt{shared\_ptr}，此时拷贝操作可以自动工作：

\begin{lstlisting}[caption={PImpl shared\_ptr 版（浅拷贝）}]
class SharedWidget {
    struct Impl;
    std::shared_ptr<Impl> pImpl_;
public:
    SharedWidget(int id);
    ~SharedWidget();   // 仍需在 .cpp 中定义

    // shared_ptr 的拷贝是浅拷贝（共享同一个 Impl）
    // 若需要深拷贝，仍须手动实现
    SharedWidget(const SharedWidget&)            = default;
    SharedWidget& operator=(const SharedWidget&) = default;
    SharedWidget(SharedWidget&&)                 = default;
    SharedWidget& operator=(SharedWidget&&)      = default;
};
\end{lstlisting}

\section{工程实战案例：PImpl 隔离重型头文件}

PImpl 最直接的工程价值在于将 ``污染性'' 头文件从公共接口中移除。以下展示三个典型的真实场景。

\subsection{案例一：隔离 WinSock 头文件}

Windows 网络编程中，\texttt{winsock2.h} 和 \texttt{windows.h} 会引入大量宏定义和类型，与用户代码产生冲突，且显著拖慢编译速度。

\subsubsection*{问题：不使用 PImpl}

\begin{lstlisting}[caption={网络类——头文件直接包含 WinSock（反面示例）}]
// tcp_client.h -- 反面示例
#pragma once
#include <winsock2.h>       // 引入 ~30,000 行宏和类型
#include <ws2tcpip.h>
#include <string>

class TcpClient {
    SOCKET sock_;            // 暴露了 Windows 平台类型
    sockaddr_in addr_;       // 暴露了 WinSock 结构体
    bool connected_;
public:
    TcpClient();
    ~TcpClient();
    bool connect(const std::string& host, int port);
    int  send(const char* data, int len);
    int  recv(char* buf, int len);
    void close();
};
\end{lstlisting}

任何包含 \texttt{tcp\_client.h} 的翻译单元都会被迫编译 WinSock 的全部声明。更严重的是，\texttt{winsock2.h} 中的 \texttt{min}/\texttt{max} 宏、\texttt{ERROR} 宏等会与标准库或用户代码冲突。

\subsubsection*{解决方案：PImpl 隔离}

\begin{lstlisting}[caption={网络类——PImpl 隔离 WinSock（头文件）}]
// tcp_client.h -- 干净的公共接口
#pragma once
#include <memory>
#include <string>
#include <span>

class TcpClient {
public:
    TcpClient();
    ~TcpClient();

    // 禁止拷贝（socket 不可复制）
    TcpClient(const TcpClient&)            = delete;
    TcpClient& operator=(const TcpClient&) = delete;

    // 允许移动
    TcpClient(TcpClient&& other) noexcept;
    TcpClient& operator=(TcpClient&& other) noexcept;

    bool connect(const std::string& host, int port);
    int  send(std::span<const char> data);
    int  recv(std::span<char> buf);
    bool isConnected() const;
    void close();

private:
    struct Impl;
    std::unique_ptr<Impl> pImpl_;
};
\end{lstlisting}

\begin{lstlisting}[caption={网络类——PImpl 隔离 WinSock（实现文件）}]
// tcp_client.cpp -- WinSock 仅在此处可见
#include "tcp_client.h"

#include <winsock2.h>     // 仅 .cpp 包含
#include <ws2tcpip.h>     // 仅 .cpp 包含
#pragma comment(lib, "ws2_32.lib")

struct TcpClient::Impl {
    SOCKET sock = INVALID_SOCKET;
    bool connected = false;

    // WinSock 初始化为进程级别操作，此处省略 RAII 封装
    ~Impl() {
        if (sock != INVALID_SOCKET)
            closesocket(sock);
    }
};

TcpClient::TcpClient()
    : pImpl_(std::make_unique<Impl>()) {}

TcpClient::~TcpClient() = default;
TcpClient::TcpClient(TcpClient&&) noexcept = default;
TcpClient& TcpClient::operator=(TcpClient&&) noexcept = default;

bool TcpClient::connect(const std::string& host, int port) {
    addrinfo hints{}, *result = nullptr;
    hints.ai_family   = AF_INET;
    hints.ai_socktype = SOCK_STREAM;
    hints.ai_protocol = IPPROTO_TCP;

    auto portStr = std::to_string(port);
    if (getaddrinfo(host.c_str(), portStr.c_str(),
                    &hints, &result) != 0)
        return false;

    pImpl_->sock = socket(result->ai_family,
                          result->ai_socktype,
                          result->ai_protocol);
    if (pImpl_->sock == INVALID_SOCKET) {
        freeaddrinfo(result);
        return false;
    }

    if (::connect(pImpl_->sock, result->ai_addr,
                  (int)result->ai_addrlen) == SOCKET_ERROR) {
        closesocket(pImpl_->sock);
        pImpl_->sock = INVALID_SOCKET;
        freeaddrinfo(result);
        return false;
    }

    freeaddrinfo(result);
    pImpl_->connected = true;
    return true;
}

int TcpClient::send(std::span<const char> data) {
    return ::send(pImpl_->sock, data.data(),
                  (int)data.size(), 0);
}

int TcpClient::recv(std::span<char> buf) {
    return ::recv(pImpl_->sock, buf.data(),
                  (int)buf.size(), 0);
}

bool TcpClient::isConnected() const {
    return pImpl_->connected;
}

void TcpClient::close() {
    if (pImpl_->sock != INVALID_SOCKET) {
        closesocket(pImpl_->sock);
        pImpl_->sock = INVALID_SOCKET;
        pImpl_->connected = false;
    }
}
\end{lstlisting}

\textbf{编译收益}：用户代码中 \texttt{\#include "tcp\_client.h"} 不再触发 WinSock 头文件的编译。在大型项目中，这可以将单个翻译单元的预处理+编译时间减少数秒，跨数百个文件累积效果显著。同时消除了宏污染问题。

\subsection{案例二：隔离重型模板库（几何内核 / Boost）}

Boost.Geometry、CGAL、Eigen 等库的头文件包含大量模板元编程代码，实例化代价极高。一个看似简单的 \texttt{\#include <boost/geometry.hpp>} 可能引入数万行模板代码。

\begin{lstlisting}[caption={几何处理器——PImpl 隔离 Boost.Geometry}]
// geometry_processor.h -- 用户接口零依赖
#pragma once
#include <memory>
#include <string>
#include <vector>

struct Point2D { double x, y; };   // 自定义轻量类型

class GeometryProcessor {
public:
    GeometryProcessor();
    ~GeometryProcessor();

    GeometryProcessor(GeometryProcessor&&) noexcept;
    GeometryProcessor& operator=(GeometryProcessor&&) noexcept;

    // 使用自定义轻量类型，不暴露 Boost 类型
    double polygonArea(const std::vector<Point2D>& ring) const;
    bool   contains(const std::vector<Point2D>& polygon,
                    Point2D point) const;
    std::vector<Point2D> convexHull(
        const std::vector<Point2D>& points) const;
    double distance(Point2D a, Point2D b) const;

private:
    struct Impl;
    std::unique_ptr<Impl> pImpl_;
};
\end{lstlisting}

\begin{lstlisting}[caption={几何处理器——实现文件隔离 Boost}]
// geometry_processor.cpp
#include "geometry_processor.h"

// 重型头文件仅在 .cpp 中可见
#include <boost/geometry.hpp>
#include <boost/geometry/geometries/polygon.hpp>
#include <boost/geometry/geometries/point_xy.hpp>
#include <boost/geometry/algorithms/area.hpp>
#include <boost/geometry/algorithms/within.hpp>
#include <boost/geometry/algorithms/convex_hull.hpp>

namespace bg = boost::geometry;
using BPoint   = bg::model::d2::point_xy<double>;
using BPolygon = bg::model::polygon<BPoint>;

struct GeometryProcessor::Impl {
    // 可缓存预处理的几何对象、空间索引等
    // 此处省略具体缓存逻辑

    static BPolygon toBoost(const std::vector<Point2D>& ring) {
        BPolygon poly;
        for (auto& p : ring)
            bg::append(poly.outer(), BPoint(p.x, p.y));
        bg::correct(poly);
        return poly;
    }
};

GeometryProcessor::GeometryProcessor()
    : pImpl_(std::make_unique<Impl>()) {}
GeometryProcessor::~GeometryProcessor() = default;
GeometryProcessor::GeometryProcessor(GeometryProcessor&&) noexcept = default;
GeometryProcessor& GeometryProcessor::operator=(
    GeometryProcessor&&) noexcept = default;

double GeometryProcessor::polygonArea(
        const std::vector<Point2D>& ring) const {
    return bg::area(pImpl_->toBoost(ring));
}

bool GeometryProcessor::contains(
        const std::vector<Point2D>& polygon,
        Point2D point) const {
    return bg::within(BPoint(point.x, point.y),
                      pImpl_->toBoost(polygon));
}

std::vector<Point2D> GeometryProcessor::convexHull(
        const std::vector<Point2D>& points) const {
    // 构建多点集 -> 计算凸包 -> 转换回自定义类型
    bg::model::multi_point<BPoint> mpts;
    for (auto& p : points)
        bg::append(mpts, BPoint(p.x, p.y));

    BPolygon hull;
    bg::convex_hull(mpts, hull);

    std::vector<Point2D> result;
    for (auto& bp : hull.outer())
        result.push_back({bg::get<0>(bp), bg::get<1>(bp)});
    return result;
}

double GeometryProcessor::distance(Point2D a, Point2D b) const {
    return bg::distance(BPoint(a.x, a.y), BPoint(b.x, b.y));
}
\end{lstlisting}

\textbf{编译收益}：Boost.Geometry 的模板实例化仅发生在 \texttt{geometry\_processor.cpp} 一个翻译单元中。所有使用 \texttt{GeometryProcessor} 的用户代码完全不需要编译 Boost 头文件。在 CI 构建中，这类隔离可以将增量编译时间从分钟级降低到秒级。

\subsection{案例三：隔离 JSON 序列化头文件}

\texttt{nlohmann/json} 是 header-only 库，单个头文件超过 25,000 行，大量模板实例化使得包含它的翻译单元编译时间急剧膨胀。

\begin{lstlisting}[caption={配置管理器——PImpl 隔离 JSON 库}]
// config_manager.h -- 不暴露 nlohmann/json
#pragma once
#include <memory>
#include <string>
#include <filesystem>
#include <optional>
#include <vector>

class ConfigManager {
public:
    explicit ConfigManager(const std::filesystem::path& filePath);
    ~ConfigManager();

    ConfigManager(ConfigManager&&) noexcept;
    ConfigManager& operator=(ConfigManager&&) noexcept;

    // 类型安全的访问接口——不暴露 json 类型
    std::optional<std::string> getString(
        const std::string& key) const;
    std::optional<int> getInt(const std::string& key) const;
    std::optional<double> getDouble(
        const std::string& key) const;
    std::optional<bool> getBool(const std::string& key) const;

    void set(const std::string& key, const std::string& value);
    void set(const std::string& key, int value);
    void set(const std::string& key, double value);
    void set(const std::string& key, bool value);

    bool save() const;
    bool reload();
    std::vector<std::string> keys() const;

private:
    struct Impl;
    std::unique_ptr<Impl> pImpl_;
};
\end{lstlisting}

\begin{lstlisting}[caption={配置管理器——实现文件隔离 JSON 库}]
// config_manager.cpp
#include "config_manager.h"

// 25,000+ 行的 header-only 库，仅此处可见
#include <nlohmann/json.hpp>
#include <fstream>

using json = nlohmann::json;

struct ConfigManager::Impl {
    std::filesystem::path filePath;
    json data;

    explicit Impl(const std::filesystem::path& path)
        : filePath(path) {
        std::ifstream ifs(path);
        if (ifs.is_open())
            data = json::parse(ifs, nullptr, false);
    }
};

ConfigManager::ConfigManager(
        const std::filesystem::path& filePath)
    : pImpl_(std::make_unique<Impl>(filePath)) {}

ConfigManager::~ConfigManager() = default;
ConfigManager::ConfigManager(ConfigManager&&) noexcept = default;
ConfigManager& ConfigManager::operator=(
    ConfigManager&&) noexcept = default;

std::optional<std::string> ConfigManager::getString(
        const std::string& key) const {
    if (auto it = pImpl_->data.find(key);
        it != pImpl_->data.end() && it->is_string())
        return it->get<std::string>();
    return std::nullopt;
}

std::optional<int> ConfigManager::getInt(
        const std::string& key) const {
    if (auto it = pImpl_->data.find(key);
        it != pImpl_->data.end() && it->is_number_integer())
        return it->get<int>();
    return std::nullopt;
}

std::optional<double> ConfigManager::getDouble(
        const std::string& key) const {
    if (auto it = pImpl_->data.find(key);
        it != pImpl_->data.end() && it->is_number())
        return it->get<double>();
    return std::nullopt;
}

std::optional<bool> ConfigManager::getBool(
        const std::string& key) const {
    if (auto it = pImpl_->data.find(key);
        it != pImpl_->data.end() && it->is_boolean())
        return it->get<bool>();
    return std::nullopt;
}

void ConfigManager::set(const std::string& key,
                        const std::string& value) {
    pImpl_->data[key] = value;
}

void ConfigManager::set(const std::string& key, int value) {
    pImpl_->data[key] = value;
}

void ConfigManager::set(const std::string& key, double value) {
    pImpl_->data[key] = value;
}

void ConfigManager::set(const std::string& key, bool value) {
    pImpl_->data[key] = value;
}

bool ConfigManager::save() const {
    std::ofstream ofs(pImpl_->filePath);
    if (!ofs.is_open()) return false;
    ofs << pImpl_->data.dump(4);
    return ofs.good();
}

bool ConfigManager::reload() {
    std::ifstream ifs(pImpl_->filePath);
    if (!ifs.is_open()) return false;
    pImpl_->data = json::parse(ifs, nullptr, false);
    return !pImpl_->data.is_discarded();
}

std::vector<std::string> ConfigManager::keys() const {
    std::vector<std::string> result;
    for (auto& [k, v] : pImpl_->data.items())
        result.push_back(k);
    return result;
}
\end{lstlisting}

\textbf{编译收益}：\texttt{nlohmann/json} 的 25,000+ 行模板代码仅在一个翻译单元中编译。项目中可能有数十个模块需要使用配置管理器，但都不需要包含 JSON 头文件。在实测中，移除一个 \texttt{nlohmann/json} 包含可以使单个翻译单元的编译时间减少 2--5 秒。

\subsection{隔离效果量化参考}

\begin{table}[ht]
\centering
\caption{PImpl 隔离重型头文件的典型编译收益}
\label{tab:isolation-benefit}
\begin{tabular}{@{}llll@{}}
\toprule
\textbf{被隔离的头文件} & \textbf{头文件规模} & \textbf{单次包含开销} & \textbf{典型影响范围} \\
\midrule
\texttt{winsock2.h + windows.h}   & $\sim$30,000 行 & 0.5--2 秒  & 网络相关全部模块 \\
\texttt{boost/geometry.hpp}       & $\sim$50,000 行 & 3--8 秒   & 几何计算全部模块 \\
\texttt{Eigen/Dense}              & $\sim$10,000 行 & 1--3 秒   & 数值计算全部模块 \\
\texttt{nlohmann/json.hpp}        & $\sim$25,000 行 & 2--5 秒   & 序列化全部模块 \\
\texttt{CGAL/Exact\_predicates.h} & $\sim$80,000 行 & 5--15 秒  & 计算几何全部模块 \\
\bottomrule
\end{tabular}
\end{table}

\noindent\textbf{注意}：表 \ref{tab:isolation-benefit} 中的数值为典型参考值，实际开销取决于编译器、优化级别和硬件配置。关键原则是：\textbf{如果一个头文件被 N 个翻译单元包含，PImpl 隔离的收益 = (单次包含开销) $\times$ (N $-$ 1)}。在大型项目中 N 往往在 50--500 之间，累积收益非常可观。

\section{总结与选择建议}

\subsection{CRTP 选型决策}

选择 CRTP 还是虚函数，核心在于\textbf{性能需求}和\textbf{运行时多态需求}的权衡：

\begin{lstlisting}[caption={决策参考：CRTP vs 虚函数 vs Deducing this}]
// ==========================================
// 场景 1：需要运行时多态（容器存放不同类型）
// -> 使用虚函数
// ==========================================
class Shape {
public:
    virtual double area() const = 0;
    virtual ~Shape() = default;
};
// std::vector<std::unique_ptr<Shape>> shapes;  // OK

// ==========================================
// 场景 2：编译期已知类型，追求零开销
// -> 使用传统 CRTP (C++11 起可用)
// ==========================================
template <typename Derived>
class ShapeCRTP {
public:
    double area() const {
        return static_cast<const Derived*>(this)->areaImpl();
    }
};
// std::vector<Circle> circles;  // 同质容器，无 vtable 开销

// ==========================================
// 场景 3：C++23 项目，追求简洁与安全
// -> 使用 Deducing this
// ==========================================
class ShapeBase {
public:
    double area(this const auto& self) {
        return self.areaImpl();  // 无需 static_cast
    }
};
\end{lstlisting}

\begin{table}[ht]
\centering
\caption{CRTP 技术选型对照}
\label{tab:crtp-decision}
\begin{tabular}{@{}llll@{}}
\toprule
\textbf{需求} & \textbf{虚函数} & \textbf{传统 CRTP} & \textbf{C++23 Deducing this} \\
\midrule
运行时多态         & 支持     & 不支持   & 不支持   \\
零运行时开销       & 否       & 是       & 是       \\
异质容器           & 支持     & 不支持   & 不支持   \\
最低 C++ 标准      & C++98    & C++98    & C++23    \\
代码简洁度         & 高       & 中       & 高       \\
类型安全           & 高       & 中       & 高       \\
编译错误可读性     & 好       & 差       & 好       \\
混入组合           & 不便     & 方便     & 方便     \\
\bottomrule
\end{tabular}
\end{table}

\subsection{PImpl 选型决策}

PImpl 的智能指针选择取决于\textbf{所有权语义}和\textbf{C++ 标准版本}：

\begin{lstlisting}[caption={决策参考：PImpl 指针类型选择}]
// ==========================================
// 场景 1：独占所有权（最常见）
// -> std::unique_ptr
// ==========================================
class Widget {
    struct Impl;
    std::unique_ptr<Impl> pImpl_;
public:
    Widget();
    ~Widget();
    Widget(const Widget&);             // 深拷贝
    Widget& operator=(const Widget&);
    Widget(Widget&&) noexcept;         // 移动
    Widget& operator=(Widget&&) noexcept;
};

// ==========================================
// 场景 2：共享所有权（多个 Widget 共享同一个 Impl）
// -> std::shared_ptr
// ==========================================
class SharedWidget {
    struct Impl;
    std::shared_ptr<Impl> pImpl_;
    // 拷贝 = 浅拷贝（共享 Impl）
    // 深拷贝需手动实现
};

// ==========================================
// 场景 3：遗留 C++03 代码，无法使用智能指针
// -> 裸指针 + 手动管理
// ==========================================
class LegacyWidget {
    struct Impl;
    Impl* pImpl_;   // 必须实现规则五
};
\end{lstlisting}

\subsection{何时不应使用这些模式}

CRTP 和 PImpl 并非银弹，以下场景应避免使用：

\begin{table}[ht]
\centering
\caption{反模式：何时不使用 CRTP / PImpl}
\label{tab:anti-pattern}
\begin{tabular}{@{}lll@{}}
\toprule
\textbf{场景} & \textbf{不推荐} & \textbf{原因与替代} \\
\midrule
需要运行时多态分发       & CRTP      & CRTP 是编译期技术，无法放入异质容器；用虚函数 \\
简单类（1--2 个成员）    & PImpl     & 过度设计，增加不必要的间接层；直接暴露成员即可 \\
对调试体验要求极高       & 传统 CRTP & 模板错误栈深，调试困难；考虑 C++23 或虚函数 \\
头文件依赖本就很少       & PImpl     & 编译防火墙收益低，反而增加 .cpp 代码量 \\
需要值语义 + 深拷贝      & shared\_ptr PImpl & shared\_ptr 默认浅拷贝，需额外手动实现 \\
嵌入式/禁用动态分配      & PImpl     & PImpl 需要堆分配；考虑预分配或栈上实现 \\
\bottomrule
\end{tabular}
\end{table}

\subsection{性能与编译时间考量}

\subsubsection{运行时性能}

\begin{itemize}[leftmargin=2em]
  \item \textbf{CRTP vs 虚函数}：CRTP 的函数调用在编译期内联，无 vtable 查找开销。在紧密循环中差异可达 10--30\%；在普通业务代码中差异可忽略。
  \item \textbf{PImpl 开销}：每次接口调用都经过指针间接寻址（一次额外的解引用）。编译器通常无法跨编译单元内联，因此热路径上的函数应考虑直接在头文件中实现。
  \item \textbf{内存开销}：PImpl 引入一次额外的堆分配（\texttt{sizeof(unique\_ptr)} + 堆上的 \texttt{Impl} 对象）。对于频繁创建销毁的小对象，可考虑对象池。
\end{itemize}

\subsubsection{编译时间}

\begin{itemize}[leftmargin=2em]
  \item \textbf{PImpl 的核心价值}：将实现细节的 \texttt{\#include} 从头文件移到 .cpp 文件，可显著减少用户代码的编译时间。在大型项目中，这可以将增量编译时间减少 50\% 以上。
  \item \textbf{CRTP 的编译开销}：模板实例化会增加编译时间，但通常比虚函数的链接开销更可控。C++23 Deducing this 的编译开销与传统 CRTP 相当。
\end{itemize}

\subsection{综合决策表}

\begin{table}[ht]
\centering
\caption{完整技术选型指南}
\label{tab:full-decision}
\begin{tabular}{@{}llll@{}}
\toprule
\textbf{场景} & \textbf{推荐技术} & \textbf{C++ 标准} & \textbf{关键理由} \\
\midrule
高性能静态多态         & 传统 CRTP           & $\geq$11  & 零成本抽象，广泛支持 \\
现代项目的静态多态     & Deducing this       & $\geq$23  & 更简洁、更安全 \\
编译防火墙             & PImpl + unique\_ptr & $\geq$11  & 标准做法，RAII 自动管理 \\
ABI 稳定的动态库接口   & PImpl + unique\_ptr & $\geq$11  & 实现变更不影响二进制布局 \\
共享资源/观察者模式    & PImpl + shared\_ptr & $\geq$11  & 浅拷贝语义，自动引用计数 \\
遗留 C++03 代码维护    & PImpl + 裸指针      & 任意      & 最小改动，无需智能指针 \\
运行时多态（异质容器） & 虚函数              & 任意      & CRTP/PImpl 均不适用 \\
混入多个行为           & CRTP 链式继承       & $\geq$11  & 可组合，避免菱形继承 \\
\bottomrule
\end{tabular}
\end{table}

\subsection{最佳实践清单}

\begin{table}[ht]
\centering
\caption{最佳实践速查}
\label{tab:best-practice}
\begin{tabular}{@{}p{0.15\textwidth}p{0.8\textwidth}@{}}
\toprule
\textbf{模式} & \textbf{要点} \\
\midrule
传统 CRTP
& 将基类析构设为 \texttt{protected}，防止通过基类指针删除派生对象；
  提供 \texttt{derived()} 辅助函数封装 \texttt{static\_cast}；
  使用 \texttt{friend} 声明限制派生类的继承关系 \\
\addlinespace
C++23 this
& 优先用 \texttt{this auto\&\&} 统一处理 cv-ref 重载；
  在递归 Lambda 中使用显式 this 替代 \texttt{std::function}；
  配合 \texttt{if constexpr} 做编译期分发 \\
\addlinespace
PImpl unique\_ptr
& 析构、移动操作必须在 .cpp 中定义（\texttt{= default}）；
  拷贝操作需手动深拷贝；
  头文件中不包含实现细节的任何 \texttt{\#include} \\
\addlinespace
PImpl shared\_ptr
& 明确语义：默认拷贝是浅拷贝；
  如需深拷贝，手动实现拷贝构造/赋值；
  注意循环引用问题，必要时使用 \texttt{weak\_ptr} \\
\bottomrule
\end{tabular}
\end{table}

\vspace{0.8cm}
\noindent\fbox{\parbox{0.95\textwidth}{\small
\textbf{一句话总结}：需要编译期多态选 CRTP（C++23 优先 Deducing this），需要编译防火墙选 PImpl + unique\_ptr，需要运行时多态选虚函数。三者可组合使用，但应避免过度设计。}}

\vspace{0.5cm}
\noindent\textcolor{gray}{\small 本文档中的代码示例仅供学习参考，实际使用时请根据项目需求和 C++ 标准版本选择合适的实现方式。}

\end{document}
